\section{Research Dossier --- Chapter: ``Dopamine, Pleasure, and the Self: The Philosophy''}\label{07_pleasure_and_self-research-dossier-chapter-dopamine-pleasure-and-the-self-the-philosophy}

\emph{(For: I Have No Dopamine, and I Must Function: A Tweaker Manifesto, Part C of the dopamine triptych)}

\subsection{TL;DR}\label{07_pleasure_and_self-tldr}

\begin{itemize}
\tightlist
\item
  The neuroscience literature provides robust ammunition for the chapter's central claim: dopamine system damage is not a moral failing but a material collapse of the substrate that makes valuing, wanting, initiating, and (in humans) experiencing pleasure possible --- a position now defensible from Volkow's PET imaging through Berridge's own concession on human studies, through neurology's ``abulia''/``non-will'' tradition, and through eliminative materialist philosophy of mind.
\item
  The Berridge/Robinson ``wanting vs.~liking'' distinction --- endlessly cited to imply that dopamine-deficient people aren't really suffering hedonically --- rests primarily on rodent orofacial taste-reactivity and is \emph{contradicted in humans} by D2 antagonist trials (Amisulpride/healthy volunteers), antipsychotic-induced anhedonia, Parkinson's depression studies, and effort-based decision-making work; the chapter can take a strong position that the wanting/liking dichotomy as applied to \emph{human phenomenology} has been smuggled in past the evidence.
\item
  The differential moralization of dopamine disorders (Parkinson's = compassion; ADHD = skepticism; addiction = punishment) is unjust on the substrate's own terms: the same imaging signature --- striatal D2/D3 hypofunction, blunted dopamine release, reduced reward-pathway engagement --- recurs across all three (Volkow et al., multiple papers). Recommend grounding the polemic in this fact and in the philosophical literature on agency-as-biological-computation (Churchland, Pereboom/Caruso), while honestly acknowledging the brain-disease model's own limitations.
\end{itemize}

\begin{center}\rule{0.5\linewidth}{0.5pt}\end{center}

\subsection{Key Findings}\label{07_pleasure_and_self-key-findings}

\subsubsection{1. RECLAIMING ``PLEASURE'' --- HISTORY AND TERMINOLOGY}\label{07_pleasure_and_self-reclaiming-pleasure-history-and-terminology}

\textbf{Cultural-religious demonization of pleasure}

\begin{itemize}
\tightlist
\item
  \textbf{Augustine of Hippo and concupiscence.} Augustine's doctrine of original sin made \emph{concupiscentia} --- disordered desire and bodily pleasure --- the vector by which sin is transmitted from Adam to all humanity. He taught that sexual pleasure is intrinsically tainted \href{https://uscatholic.org/articles/202411/what-did-st-augustine-say-about-original-sin/}{U.S. Catholic} (``the act itself cannot be effected without ardour of lust'') \href{https://www.arcaneknowledge.org/catholic/original5.htm}{Arcaneknowledge} and that even legitimate marital intercourse carries venial sin from this source. \href{https://www.arcaneknowledge.org/catholic/original5.htm}{Arcaneknowledge} The \emph{Confessions} and \emph{De nuptiis et concupiscentia} are the locus classicus. This installed pleasure-as-defect at the center of Western moral psychology for \textasciitilde1,500 years. (Sources: Wikipedia, ``Original sin''; ``Concupiscence''; ``Augustine of Hippo.'' See also Brachtendorf and Trapè scholarship cited there.)
\item
  \textbf{Puritanism and Calvinist asceticism.} Despite revisionist defenders (Doriani, \emph{Westminster Theological Journal} 53.1, 1991; Packer's ``Christian hedonism'' school), the historical record documents Puritan preachers (Baxter, Henry) treating sensory pleasure as fundamentally dangerous: ``Strive more against your own flesh, than against all your Enemies in Earth and Hell'' \href{https://www.desiringgod.org/messages/christian-hedonists-or-religious-prudes-the-puritans-on-sex}{Desiring God} (Baxter). Max Weber's \emph{Protestant Ethic and the Spirit of Capitalism} identified Puritan ``inner-worldly asceticism'' \href{https://www.litcharts.com/lit/the-protestant-ethic-and-the-spirit-of-capitalism/themes/puritan-asceticism}{LitCharts} as a third pillar of the capitalist ethos --- work as a hedge against pleasure.
\item
  \textbf{The deep structure} --- visible from Plato's \emph{Phaedo} through Pauline epistles through monasticism, Jansenism, Methodism, and Victorian respectability --- is that pleasure is at best a permitted concession and at worst the wellspring of evil. This is the cultural water in which the neuroscience of reward swims.
\end{itemize}

\textbf{Olds and Milner and the ``pleasure center'' --- and its retreat}

\begin{itemize}
\tightlist
\item
  Olds and Milner (1954, \emph{Journal of Comparative and Physiological Psychology} 47:419--427) discovered that rats would self-stimulate electrodes in the septal area / medial forebrain bundle until exhaustion, identifying what was immediately called the brain's ``pleasure center.'' The original framing was unabashedly hedonic: reinforcement = reward = pleasure.
\item
  Robert Heath's notorious B-19 case (1972) cemented the popular image: subjects pressed buttons ``thousands of times,'' reporting ``feelings of pleasure, alertness, and warmth'' \href{https://en.wikipedia.org/wiki/Brain_stimulation_reward}{Wikipedia} and pleading not to have the unit removed.
\item
  \textbf{The terminological retreat} is well documented and traceable: by the 1980s ``pleasure center'' had been replaced in the literature by ``reward pathway'' (more neutral, behaviorist); by the 1990s Berridge and Robinson had introduced ``incentive salience'' and ``wanting''; by the 2000s Berridge \& Kringelbach (Neuron 2015 review, \emph{Pleasure Systems in the Brain}) were explicitly stating: ``some of the best known textbook candidates for pleasure generators, including classic pleasure electrodes and the mesolimbic dopamine system, may not generate pleasure after all.'' Kringelbach \& Berridge (2017): ``in contrast to early views of dopamine as a pleasure neurotransmitter---this research has shown that dopamine is not linked to pleasure per se but instead to incentive salience.'' \href{https://sites.lsa.umich.edu/berridge-lab/wp-content/uploads/sites/743/2019/10/2017-Kringelbach-Berridge-Affective-core-Emotion-Rev-1.pdf}{University of Michigan}
\item
  \textbf{Scholarship on cultural forces shaping this terminology.} Scott Vrecko has done the most explicit historical-sociological work: ``Birth of a brain disease: science, the state and addiction neuropolitics'' (\emph{History of the Human Sciences} 23, 2010, 52--67) and ``\,`Civilizing technologies' and the control of deviance'' (\emph{BioSocieties} 5(1), 2010, 36--51). \href{https://www.semanticscholar.org/paper/\%E2\%80\%98Civilizing-technologies\%E2\%80\%99-and-the-control-of-Vrecko/658114856559331138582d44432532d544426fc2}{Semantic Scholar} Vrecko argues that the contemporary neurosciences of reward function as ``civilizing technologies'' --- re-coding moral problems about pleasure and self-control into bioscientific problems about cravings, but preserving the cultural project of regulating desire. \href{https://kclpure.kcl.ac.uk/portal/en/publications/civilizing-technologies-and-the-control-of-deviance(b0f2de9e-ab33-4053-8520-a838b6acbdf6).html}{King's College London} The special issue of \emph{History of the Human Sciences} (2010, ``Neuroscience, Power and Culture'') edited by Vrecko is the best entry point. \href{https://journals.sagepub.com/doi/10.1177/0952695109354395}{Sage Journals} \href{http://somatosphere.net/2010/neuroscience-power-and-culture-special.html/}{Somatosphere} See also Hammer et al., ``Addiction: Current Criticism of the Brain Disease Paradigm'' (PMC3969751), which documents how addiction terminology (``dopamine head,'' ``subcortical lives'') \href{https://pmc.ncbi.nlm.nih.gov/articles/PMC3969751/}{PubMed Central} imports moralized imagery.
\end{itemize}

\textbf{Argument that dopamine IS, functionally, the pleasure molecule for human experience}

The author's case can be built honestly: even granting Berridge's rodent dissociation, human evidence shows dopamine is necessary for hedonic experience as humans actually live it.

\begin{itemize}
\tightlist
\item
  \textbf{D2 blockade impairs hedonic responsivity in humans.} Weber et al., ``Does partial blockade of dopamine D2 receptors with Amisulpride cause anhedonia? An experimental study in healthy volunteers'' (\emph{Journal of Psychiatric Research}, 2023; PubMed 36680855): single 300 mg dose of amisulpride vs.~placebo, n=85, double-blind RCT. \href{https://pubmed.ncbi.nlm.nih.gov/36680855/}{PubMed} Result: significantly lower positivity ratings across all stimulus categories (p=.026, f=0.25) and lower electrodermal responses (p=.017, f=0.27). \href{https://pubmed.ncbi.nlm.nih.gov/36680855/}{PubMed} Authors' conclusion: ``D2 blockade via Amisulpride can reduce at-the-moment hedonic responsivity in healthy volunteers\ldots{} antipsychotic medication contributes to clinical anhedonia, probably via antagonistic effects at the dopamine D2 receptor.'' \href{https://pubmed.ncbi.nlm.nih.gov/36680855/}{PubMed}
\item
  \textbf{Antipsychotic-induced anhedonia.} First-generation antipsychotics produce ``loss of drive, energy, and motivation; apathy, anhedonia''; \href{https://pmc.ncbi.nlm.nih.gov/articles/PMC4826766/}{PubMed Central} clinical depression in 20--40\% of patients on FGAs \href{https://pmc.ncbi.nlm.nih.gov/articles/PMC4826766/}{PubMed Central} (PMC4826766, ``Inhibition of the reward system by antipsychotic treatment''). Patient self-reports describe ``a distinctive experience characterized by sedation, cognitive slowing, emotional blunting and reduced motivation'' \href{https://en.wikipedia.org/wiki/Anhedonia}{Wikipedia} \href{https://www.sciencedirect.com/science/article/abs/pii/S0022395623000146}{ScienceDirect} (Thompson et al.~2020).
\item
  \textbf{Parkinson's anhedonia is real.} Costello et al., ``Impaired reward sensitivity in Parkinson's depression is unresponsive to dopamine treatment'' (\emph{Brain}, 2025; PMC12129732): PD depression ``characterized by lower acceptance of offers, driven by markedly lower incentivization by reward (reward sensitivity), compared to all other groups.'' \href{https://www.ncbi.nlm.nih.gov/pmc/articles/PMC12129732/}{nih} Music-evoked pleasure: Cameron et al.~(bioRxiv 2023) found PD patients show flattened ``pleasurable urge to move to music,'' directly linking dopamine to musical pleasure.
\item
  \textbf{Pramipexole (D2/D3 agonist) modulates taste hedonics in healthy volunteers} (PMC9515131): improved taste recognition, altered hedonic evaluation. \href{https://www.ncbi.nlm.nih.gov/pmc/articles/PMC9515131/}{nih}
\item
  \textbf{Berridge himself acknowledges the gap with humans.} From the 2016 review (PMC5171207): Leyton et al.~(2005) and Brauer \& De Wit (1997) found dopamine suppression ``did not reduce pleasure ratings of drug rewards''; \href{https://pmc.ncbi.nlm.nih.gov/articles/PMC5171207/}{PubMed Central} but Hardman et al.~2012 and Sienkiewicz-Jarosz et al.~2013 \href{https://pmc.ncbi.nlm.nih.gov/articles/PMC5171207/}{PubMed Central} are typically cited \emph{for} the dissociation in humans. \href{https://pmc.ncbi.nlm.nih.gov/articles/PMC5171207/}{PubMed Central} The chapter can fairly point out: the human evidence is \emph{much weaker and more equivocal} than the rodent evidence, and the strongest \emph{experimental} human data (Weber et al.~2023, Costello et al.~2025) point the \emph{other} way.
\end{itemize}

\begin{center}\rule{0.5\linewidth}{0.5pt}\end{center}

\subsubsection{2. CRITIQUE OF WANTING/LIKING AS APPLIED TO HUMANS}\label{07_pleasure_and_self-critique-of-wantingliking-as-applied-to-humans}

\textbf{The framework's claims}

\begin{itemize}
\tightlist
\item
  Origin: Berridge, Venier \& Robinson (1989, \emph{Behavioral Neuroscience} 105:3--14): 6-OHDA dopamine lesions abolished motivation in rats but preserved orofacial ``liking'' reactions to sucrose.
\item
  Major statements: Berridge \& Robinson (1998, \emph{Brain Research Reviews} 28:309--369, ``What is the role of dopamine in reward?''); \href{https://www.sciencedirect.com/science/article/abs/pii/104457659290011P}{ScienceDirect} Robinson \& Berridge (1993 \emph{Brain Res Rev} 18:247--291, IST original); Berridge, Robinson \& Aldridge (2009 \emph{Curr Opin Pharmacol} 9:65--73); Berridge \& Kringelbach (2015 \emph{Neuron} ``Pleasure Systems in the Brain''); Berridge \& Robinson (2016 \emph{American Psychologist} 71:670--679); Robinson \& Berridge (2025 \emph{Annual Review of Psychology} 76:29--58 --- 30-years-on review).
\item
  Core claim: dopamine = ``wanting'' (incentive salience); \href{https://pmc.ncbi.nlm.nih.gov/articles/PMC3274778/}{PubMed Central} opioids in small subcortical ``hedonic hotspots'' (nucleus accumbens shell, ventral pallidum, parabrachial nucleus) = ``liking.'' \href{https://www.researchgate.net/publication/5539913_Berridge_KC_Kringelbach_ML_Affective_neuroscience_of_pleasure_reward_in_humans_and_animals_Psychopharmacology_Berl_199_457-480}{ResearchGate}
\end{itemize}

\textbf{The actual neurobiology --- what the rodent data show vs.~what gets claimed for humans}

\begin{itemize}
\tightlist
\item
  The hedonic hotspot data are \emph{rodent orofacial taste reactivity} --- gapes vs.~tongue protrusions to sucrose/quinine (Berridge \& Grill, derived from Steiner 1973). The ``liking'' measure is a behavioral facial proxy, not a phenomenological report.
\item
  Ventral pallidum hotspot: a small lesion can flip sucrose's hedonic valence to disgust \href{https://pmc.ncbi.nlm.nih.gov/articles/PMC5171207/}{PubMed Central} (Berridge \& Kringelbach 2015; Ho \& Berridge 2014). This is striking but speaks to opioid/GABAergic \emph{amplification} of liking, not to whether dopamine is unnecessary for pleasure as humans experience it.
\item
  The dopamine ``wanting'' data are largely from amphetamine sensitization in rats and 6-OHDA lesions producing aphagia paradoxes.
\end{itemize}

\textbf{Critical analysis: extension to humans is fragile}

\begin{itemize}
\tightlist
\item
  \textbf{Near-total absence of causal primate evidence for the wanting/liking dissociation.} The classic Berridge claim that the dissociation generalizes to primates rests on extrapolation. Wang et al.~(PMC4600774, ``Severe dopaminergic neuron loss in rhesus monkey brain impairs morphine-induced conditioned place preference'') \href{https://www.ncbi.nlm.nih.gov/pmc/articles/PMC4600774/}{nih} shows the opposite --- heavy DA loss in rhesus impairs reward behavior writ large. There is no monkey study isolating an orofacial-style ``liking'' signal that survives DA depletion in the way rodent studies claim.
\item
  \textbf{Human experimental dissociation evidence is weak.} Brauer \& De Wit (1997), Leyton et al.~(2005, cocaine), Wachtel et al.~(2002, amphetamine) --- the ``go-to'' human studies --- used acute phenylalanine/tyrosine depletion (APTD), \href{https://robinsonlab.research.wesleyan.edu/files/2014/01/Robinson-2015-Curr-Top-Behav-Neurosci.pdf}{Wesleyan} which produces only modest dopamine reduction, and depended on subjective ``high'' ratings that have face-validity problems. Effect sizes are small; replication across paradigms is inconsistent.
\item
  \textbf{Strong contrary evidence in humans.} Already enumerated (§1): amisulpride reduces healthy hedonic responsivity; antipsychotics induce anhedonia clinically; Parkinson's depression shows anhedonia \emph{unresponsive} even to dopamine treatment in some forms (Costello 2025); \href{https://www.ncbi.nlm.nih.gov/pmc/articles/PMC12129732/}{nih} ADHD imaging links D2/D3 hypofunction to \emph{both} motivation deficits \emph{and} drug ``liking'' responses (Volkow et al.~2007a, 2009); pramipexole alters hedonic taste evaluation.
\item
  \textbf{The PD ``anhedonia paradox.''} Berridge cites Sienkiewicz-Jarosz et al.~(2013) --- ``sweet liking in PD patients'' \href{https://pmc.ncbi.nlm.nih.gov/articles/PMC5171207/}{PubMed Central} --- to argue PD patients still ``like.'' \href{https://link.springer.com/chapter/10.1007/978-3-031-10477-0_17}{Springer} But the same literature documents that PD patients can no longer \emph{anticipate}, \emph{initiate}, or \emph{be moved by} what they nominally rate as pleasant \href{https://www.frontiersin.org/journals/psychiatry/articles/10.3389/fpsyt.2021.766692/full}{Frontiers} --- which is the part of pleasure that matters phenomenologically.
\item
  \textbf{Salamone's effort-based critique.} Salamone \& Correa (2002, 2012, \emph{Neuron} 76:470--485, ``The mysterious motivational functions of mesolimbic dopamine''; 2024 \emph{Annual Review of Psychology} 75:1--32) \href{https://www.annualreviews.org/content/journals/10.1146/annurev-psych-020223-012208}{Annual Reviews} reframes the data: dopamine is not specifically about ``wanting'' vs.~``liking'' but about \emph{activational} aspects of motivation --- overcoming effort costs, vigor, persistence. \href{https://pmc.ncbi.nlm.nih.gov/articles/PMC2817983/}{PubMed Central} \href{https://www.annualreviews.org/content/journals/10.1146/annurev-psych-020223-012208}{Annual Reviews} This subtly undermines the cleanness of the Berridge dichotomy and is more compatible with the lived experience of dopamine deficiency: not ``I want it but don't like it'' but ``I cannot summon the activational substrate to engage.''
\item
  \textbf{Anselme \& Robinson's own caveat (2016 paper, \emph{J Exp Psychol Anim Learn Cogn}).} Even the Berridge camp now distinguishes ``wanting''/``liking'' as \emph{unconscious} motivational/hedonic processes from \emph{conscious} wanting and liking \href{https://www.apa.org/pubs/highlights/spotlight/issue-69}{APA} --- which collapses much of the distinction's usefulness for adjudicating phenomenological claims about humans.
\end{itemize}

\textbf{How wanting/liking is weaponized to minimize suffering}

\begin{itemize}
\tightlist
\item
  \textbf{Public-facing framing.} The APA spotlight on Anselme \& Robinson (2016) --- ``the dissociation between `wanting' and `liking' can be puzzling: why would someone want a reward if they do not like it?'' \href{https://www.apa.org/pubs/highlights/spotlight/issue-69}{APA} --- is widely echoed in popular accounts (\emph{Scientific American}: ``Dopamine: The Currency of Desire''; ``ADHD and the Motivation That Never Comes,'' \emph{Psychology Today} 2026). The implicit message is: addicts don't even \emph{enjoy} their drugs, so the suffering is a kind of category error, \href{https://www.apa.org/pubs/highlights/spotlight/issue-69}{APA} and ADHD/depression ``anhedonia'' is more like a wanting deficit than a real loss of pleasure. This is exactly the move the chapter can indict.
\item
  \textbf{The chapter's claim is defensible:} the wanting/liking distinction has been deployed to naturalize an epistemic gerrymander in which whatever dopamine-deficient people \emph{report} about their inner lives is reinterpreted as a malfunction of ``wanting'' rather than as the real loss of meaningful experience that it is.
\end{itemize}

\textbf{Useable quotes for the chapter}

\begin{itemize}
\tightlist
\item
  Berridge \& Kringelbach (2015, Neuron): ``some of the best known textbook candidates for pleasure generators, including classic pleasure electrodes and the mesolimbic dopamine system, may not generate pleasure after all.''
\item
  Weber et al.~(2023): ``D2 blockade via Amisulpride can reduce at-the-moment hedonic responsivity in healthy volunteers\ldots{} antipsychotic medication contributes to clinical anhedonia.'' \href{https://pubmed.ncbi.nlm.nih.gov/36680855/}{PubMed}
\item
  Costello et al.~(2025): PD depression is characterized by ``markedly lower incentivization by reward\ldots{} this therapeutic effect {[}of dopamine{]} is not present in PD patients with depression.'' \href{https://www.ncbi.nlm.nih.gov/pmc/articles/PMC12129732/}{nih}
\end{itemize}

\begin{center}\rule{0.5\linewidth}{0.5pt}\end{center}

\subsubsection{3. ``WILLPOWER'' AS FOLK-PSYCHOLOGICAL ILLUSION}\label{07_pleasure_and_self-willpower-as-folk-psychological-illusion}

\textbf{Eliminative materialism --- Paul and Patricia Churchland}

\begin{itemize}
\tightlist
\item
  Paul Churchland, ``Eliminative Materialism and the Propositional Attitudes,'' \emph{Journal of Philosophy} 78(2), 1981, 67--90. Argues that folk psychology is a \emph{theory} --- and a bad one: ``First, the eliminative materialist will point to the widespread explanatory, predictive, and manipulative failures of folk psychology. So much of what is central and familiar to us remains a complete mystery from within folk psychology.'' \href{https://danielwharris.com/teaching/101/readings/Churchland.pdf}{Danielwharris} The fate of ``belief, desire, fear, sensation, pain, joy, and so on'' \href{https://danielwharris.com/teaching/101/readings/Churchland.pdf}{Danielwharris} is to go the way of phlogiston, witches, and demonic possession once neuroscience matures. Useful direct quote: ``Our explanations of one another's behavior will appeal to such things as our neuropharmacological states, the neural activity in specialized anatomical areas, and whatever other states are deemed relevant by the new theory.'' \href{https://danielwharris.com/teaching/101/readings/Churchland.pdf}{Danielwharris}
\item
  Patricia Churchland, \emph{Neurophilosophy} (MIT Press, 1986); \emph{Brain-Wise: Studies in Neurophilosophy} (MIT, 2002); \emph{Touching a Nerve: The Self as Brain} (Norton, 2013). The 2013 book contains a chapter explicitly titled ``Free Will, Habits, and Self-Control'' \href{https://franklin.library.upenn.edu/catalog/FRANKLIN_9960549203503681}{Franklin} --- directly relevant. P. Churchland's central move is to ground ``self,'' ``decision,'' and ``will'' in biology: ``What happens when we accept that everything we feel and think stems not from an immaterial spirit but from electrical and chemical activity in our brains?'' \href{https://books.google.com/books/about/Touching_a_Nerve.html?id=5QQSAAAAQBAJ}{Google Books}
\item
  Stanford Encyclopedia of Philosophy entry by William Ramsey (revised 2019) for cite-checkable summary.
\item
  \textbf{Honest caveat to flag:} strong eliminativism is a minority position; even some materialists (Bickle, Dennett) reject full elimination of propositional attitudes. The chapter need not commit to \emph{full} elimination --- the weaker, sufficient claim is that ``willpower'' specifically is a poor folk construct that does not pick out a real biological kind.
\end{itemize}

\textbf{Putnam, multiple realizability, and what it means here}

\begin{itemize}
\tightlist
\item
  Hilary Putnam (1967, ``Psychological Predicates''; reprinted as ``The Nature of Mental States''): mental kinds are multiply realizable --- i.e., the same mental state can be implemented in different physical substrates (humans, octopi, hypothetical Martians, silicon). This was originally an \emph{anti}-reductionist argument. \textbf{Care needed:} the author should not over-cite Putnam as a friend of the polemic. Putnam's argument is consistent with the chapter's central claim only on this reading: ``willpower,'' ``agency,'' ``wanting'' are functional/computational properties, and in humans those functions are realized through dopamine-modulated cortico-striato-thalamo-cortical loops. Damaging the realization damages the function. The chapter wants the implementation point, not the autonomy-from-physics point.
\item
  Putnam himself moved away from functionalism (1988, \href{https://en.wikipedia.org/wiki/Hilary_Putnam}{Wikipedia} \emph{Representation and Reality}).
\item
  The cleanest move: cite Putnam (1967) for the framework that mental states are functional/computational; cite Churchland for the subsequent claim that, in \emph{humans}, these functions are \emph{implemented} through specific neurochemistry, and that damage to the implementation produces real loss.
\end{itemize}

\textbf{Connection to dopamine science}

The chapter can lay out the implementation argument precisely:

\begin{itemize}
\tightlist
\item
  \textbf{PFC goal-maintenance} depends on D1-receptor-mediated tuning (Williams \& Goldman-Rakic 1995; Arnsten 2011; Cools \& D'Esposito 2011 --- the inverted-U).
\item
  \textbf{Striatal direct/indirect (Go/NoGo) pathways} --- D1-expressing direct pathway and D2-expressing indirect pathway \href{https://www.biorxiv.org/content/10.1101/2022.10.29.514339.full.pdf}{biorxiv} (Frank, Seeberger \& O'Reilly 2004 \emph{Science}; Kravitz et al.~2010 \emph{Nature}). Action selection IS this circuit.
\item
  \textbf{Limbic valuation} (NAc, OFC, vmPFC) is dopamine-modulated (Schultz 1998, 2007; Haber \& Knutson 2010).
\item
  \textbf{Thalamic gating} of cortical access depends on basal ganglia output, itself dopaminergically tuned.
\item
  Volkow's ``Disease of Free Will'' framing makes this case from inside NIDA: chronic drug use ``leads to down-regulation of prefrontal cortical activity mediated by dopamine, which in turn disrupts human capacities to `\ldots resist immediate rewards. To delay gratification. To be able to consider a goal for our future, and to carry it through.'\,'' \href{https://www.sciencedirect.com/science/article/abs/pii/S0306460321001404}{ScienceDirect} (Volkow 2015 APA address; quoted in Henne et al., ``Effects of addiction science on conceived freewill and responsibility,'' ScienceDirect S0306460321001404.)
\end{itemize}

\textbf{Free will, determinism, neuroscience}

\begin{itemize}
\tightlist
\item
  Derk Pereboom, \emph{Free Will, Agency, and Meaning in Life} (Oxford, 2014); \emph{Living Without Free Will} (Cambridge, 2001). \href{https://philpapers.org/rec/CARPOD-2}{PhilPapers} Hard incompatibilist.
\item
  Gregg Caruso, \emph{Rejecting Retributivism: Free Will, Punishment, and Criminal Justice} (Cambridge, 2021). \href{https://www.cambridge.org/core/books/moral-responsibility-reconsidered/0DD200C8E84E6458EA56A098BC7E206F}{Cambridge Core} Free will skepticism applied to criminal justice.
\item
  Pereboom \& Caruso, ``Hard-Incompatibilist Existentialism: Neuroscience, Punishment, and Meaning in Life'' in Caruso \& Flanagan (eds.), \emph{Neuroexistentialism} (Oxford, 2018). \href{https://philpapers.org/rec/CARPOD-2}{PhilPapers} \href{https://cornell.academia.edu/DerkPereboom}{Academia} The \emph{Neuroexistentialism} volume itself is the perfect anchor for the chapter's existential framing.
\item
  These thinkers are the right citations for: agency-as-biological-computation does not preclude meaning, but \emph{does} preclude the retributive-blame architecture that underwrites differential moralization.
\end{itemize}

\textbf{The ``willpower'' myth specifically}

\begin{itemize}
\tightlist
\item
  Ego depletion --- the Baumeister/Tice ``willpower as glucose-fueled muscle'' model (Baumeister, Vohs \& Tice 2007) --- has been substantially undermined by replication failures (Hagger et al.~2016 multi-lab; Vohs et al.~2021 RRR; meta-analyses post-2016). Inzlicht and colleagues have argued that high self-control is a function of \emph{not experiencing temptation as struggle} in the first place \href{https://www.goodlifeproject.com/podcast/the-surprising-truth-about-willpower-why-self-control-may-not-be-what-you-think-michael-inzlicht-phd/}{Good Life Project} --- which is an argument that ``willpower'' as folk-construct mis-describes what actually goes on in successful self-regulators.
\end{itemize}

\begin{center}\rule{0.5\linewidth}{0.5pt}\end{center}

\subsubsection{4. DIFFERENTIAL MORALIZATION OF DOPAMINE DISORDERS}\label{07_pleasure_and_self-differential-moralization-of-dopamine-disorders}

\textbf{The substrate is shared}

The author's argument that the same dopamine pathology underlies disorders we treat morally very differently is well supported:

\begin{itemize}
\tightlist
\item
  \textbf{Addiction:} Volkow et al.~(multiple, 1990s--2010s) --- chronic stimulant users show reduced striatal D2/D3 receptor availability (raclopride PET) and blunted methylphenidate-induced dopamine release. \href{https://psychiatryonline.org/doi/10.1176/pn.39.11.0390032}{Psychiatric News} Methamphetamine users show reduced dopamine transporter availability lasting years (Chou 2007; McCann 1998, 2008; Sekine 2001; Volkow 2001b). PMC4106030 (Trifilieff \& Martinez, ``Imaging addiction'').
\item
  \textbf{Obesity:} Wang, Volkow et al.~--- morbidly obese subjects \href{https://pubmed.ncbi.nlm.nih.gov/18640912/}{PubMed} show reduced striatal D2 receptor availability, comparable to addicted subjects (Wang et al.~2001; Volkow \& Wise 2005 \emph{Nat Neurosci}). PMC2696819, PubMed 18640912.
\item
  \textbf{ADHD:} Volkow et al.~2007a, 2007b, 2007c --- drug-naïve ADHD adults show lower D2/D3 receptor availability and lower DA transporters in midbrain and nucleus accumbens. \href{https://pmc.ncbi.nlm.nih.gov/articles/PMC3010326/}{PubMed Central} ``Motivation Deficit in ADHD is Associated with Dysfunction of the Dopamine Reward Pathway'' (PMC3010326): Achievement scale lower in ADHD (11±5 vs 14±3, p\textless0.001), correlated with NAc D2/D3 (r=0.39, p\textless0.008). \href{https://pmc.ncbi.nlm.nih.gov/articles/PMC3010326/}{PubMed Central}
\item
  \textbf{Parkinson's:} Substantia nigra dopaminergic cell loss; apathy in 13.9--70\% of patients (depending on definition), often \emph{preceding} motor symptoms by years (PMC8775593, PMC9313138). Apathy is now recognized as the ``Park Apathy'' \href{https://pmc.ncbi.nlm.nih.gov/articles/PMC9313138/}{PubMed Central} non-motor subtype.
\item
  \textbf{Schizophrenia negative symptoms:} Hypodopaminergia in mesocortical projections; antipsychotics (D2 antagonists) often worsen apathy/anhedonia secondarily.
\end{itemize}

\textbf{The injustice of differential moralization}

\begin{itemize}
\tightlist
\item
  \textbf{Parkinson's = compassion.} ``Apathy can be misinterpreted as laziness, poor initiative or depression,'' Michael J. Fox Foundation explicitly notes --- but in the PD context, no one is blaming patients. APDA, Parkinson Society BC, MJFF all explicitly frame apathy as \emph{biological}, not character-based.
\item
  \textbf{ADHD = skepticism.} ADHD remains widely doubted as a ``real'' disorder despite the imaging evidence above. The cultural framing (``just get organized,'' ``kids these days,'' ``everyone has trouble focusing'') tracks neither the imaging nor the clinical reality.
\item
  \textbf{Addiction = punishment.} Despite decades of NIDA/Volkow advocacy for the brain-disease model, U.S. drug policy remains overwhelmingly carceral. Stigma persists at high levels in both lay and clinical populations (PMC10867215 multinational review). \href{https://www.ncbi.nlm.nih.gov/pmc/articles/PMC10867215/}{nih}
\end{itemize}

\textbf{The ethical literature}

\begin{itemize}
\tightlist
\item
  Volkow's ``Addiction Is a Disease of Free Will'' (NIDA, 2015) is the strongest in-house statement linking dopamine deficit to diminished moral responsibility. Useful quote: ``We can do much to reduce the shame and the stigma of drug addiction, once medical professionals, and we as a society, understand that addiction is not just `a disease of the brain,' but one in which the circuits that enable us to exert free will no longer function as they should.'' \href{https://www.breakingthecycles.com/2015/06/21/disease-of-addiction-disease-of-free-will-dr-nora-volkow/}{Breakingthecycles}
\item
  Heyman, Lewis (Marc Lewis, \emph{The Biology of Desire}, 2015), and Heather (PMC5486515, ``Q: Is Addiction a Brain Disease or a Moral Failing? A: Neither'') provide the major \emph{critiques} of the brain-disease model --- the chapter should engage these, not duck them, since the polemic is stronger if it can absorb the criticism that the disease label can also stigmatize.
\item
  Hanna Pickard's \emph{responsibility-without-blame} framework (Neuroethics 10(1):169--180, 2017; \href{https://link.springer.com/article/10.1007/s12152-016-9295-2}{Springer} \emph{A Philosophy of Addiction}, Princeton 2026) \href{https://philosophy.jhu.edu/directory/hanna-pickard/}{Philosophy} is the most sophisticated current proposal --- preserves agency while removing blame. \href{https://philpapers.org/rec/PICRWB-4}{PhilPapers} The chapter can position itself relative to Pickard: agreeing on the rejection of blame; pushing harder on the substrate's reality.
\item
  \textbf{Honest caveat:} The brain-disease model has been criticized empirically (Heilig 2021 vs.~Heather et al.~2022; \emph{The Lancet Psychiatry} 2025 ``Reevaluating the brain disease model''). The chapter doesn't have to swallow the disease label whole; it can argue specifically that \emph{dopamine system damage is a material fact} without committing to all the pragmatic baggage of ``BDMA.'' \href{https://www.ncbi.nlm.nih.gov/pmc/articles/PMC5486515/}{nih}
\item
  \textbf{Buchman 2010} (``addiction-as-disease debunks the moralized argument that addiction is a problem for weak-willed people'' \href{https://pmc.ncbi.nlm.nih.gov/articles/PMC3969751/}{PubMed Central} --- see Hammer et al.~PMC3969751) is a useful framing source.
\end{itemize}

\begin{center}\rule{0.5\linewidth}{0.5pt}\end{center}

\subsubsection{5. EXISTENTIAL STAKES}\label{07_pleasure_and_self-existential-stakes}

\textbf{Clinical: the vocabulary of dopamine collapse}

The neurology and psychiatry literatures already speak the language the chapter wants to claim philosophically.

\begin{itemize}
\tightlist
\item
  \textbf{Abulia / aboulia} --- from Greek \emph{aboulia}, ``non-will.'' \href{https://www.ncbi.nlm.nih.gov/sites/books/NBK537093/}{NCBI} StatPearls (NBK537093): ``a syndrome of hypofunction\ldots{} lack of initiative, spontaneity, and drive, apathy, slowness of thought (bradyphrenia), and blunting of emotional responses.'' \href{https://www.ncbi.nlm.nih.gov/sites/books/NBK537093/}{NCBI} First described 1838. \href{https://www.ncbi.nlm.nih.gov/sites/books/NBK537093/}{NCBI} Mechanism: ``dopamine-dependent circuitry malfunction. Lesions anywhere in the `centro-medial core' of brain frontal-subcortical circuitry, from frontal lobes to the brainstem, may produce this condition.'' \href{https://www.ncbi.nlm.nih.gov/sites/books/NBK537093/}{NCBI} Treatment: levodopa, bromocriptine, bupropion (all dopaminergic). \href{https://mdsearchlight.com/neurology/abulia/}{MD Searchlight} The medical name for what the chapter is describing IS \emph{non-will}.
\item
  \textbf{Akinetic mutism, athymia, ``loss of psychic self-activation''} --- the spectrum.
\item
  \textbf{Apathy in PD:} anhedonia, loss of ``get up and go,'' loss of ``interest in people or activities, and sex'' \href{https://www.ncbi.nlm.nih.gov/pmc/articles/PMC8775593/}{nih} (PMC8775593 cohort study, n=120). MJFF: ``Hobbies and social activities may no longer bring enjoyment, and daily routines may require more energy. Basic tasks may be difficult to start and complete.'' \href{https://www.michaeljfox.org/news/ask-md-apathy-and-parkinsons-disease}{Michael J. Fox Foundation}
\end{itemize}

\textbf{The Japanese ``remarkable disturbance of spontaneity''}

The author's specific reference is most likely to one or both of these traditions --- recommend verifying the exact source the author has in mind:

\begin{itemize}
\tightlist
\item
  \textbf{Bleuler's ``disturbance of volition''} --- already in 1911 --- was the foundational Western description: ``weakening of those emotional activities which permanently form the mainspring of volition\ldots{} loss of mastery over volition, of endeavour, and of ability for independent action.'' (Quoted in PMC2833114, Foussias \& Remington, ``Negative Symptoms in Schizophrenia: Avolition and Occam's Razor.'') ``Lack of spontaneity'' (N6) and ``disturbance of volition'' (G13) are explicit PANSS items. \href{https://pmc.ncbi.nlm.nih.gov/articles/PMC2833114/}{PubMed Central}
\item
  \textbf{Kasahara Yomishi}, Kyoto/Nagoya University, 1971 onward: theorized ``student apathy'' \href{https://pmc.ncbi.nlm.nih.gov/articles/PMC11114399/}{PubMed Central} (taigakushō) as a distinct syndrome (Kasahara 1971, \href{https://www.ovid.com/journals/pacnr/fulltext/10.1002/pcn5.120~adolescent-psychiatry-of-y-kasahara-and-succeeding-research}{Ovid} 1984 \emph{Apathy syndrome}, Iwanami; \href{https://journals.sagepub.com/doi/10.1177/070674379403901012}{Sage Journals} PMC11114399, ``Adolescent psychiatry of Y. Kasahara''). The link from Kasahara to \emph{hikikomori} is well-traced. \href{https://pmc.ncbi.nlm.nih.gov/articles/PMC11114399/}{PubMed Central}
\item
  \textbf{Shimazono Yutaka} (1974, ``Prognosis of the functional psychosis,'' in Mitsuda \& Fukuda eds., \emph{Biological Mechanisms of Schizophrenia}, Tokyo: Igaku Shoin) \href{https://www.cambridge.org/core/journals/the-british-journal-of-psychiatry/article/abs/longterm-followup-study-of-schizophrenia-in-japan-with-special-reference-to-the-course-of-social-adjustment/D2214EA6A260E08FE813B31AC2116448}{Cambridge Core} is also frequently cited in long-term outcome studies of schizophrenia in Japan.
\item
  Japanese psychiatric terminology like \emph{jinkaku suijun no teika} (``decline in the level of personality'') explicitly describes negative symptoms in volitional/character terms (BJPsych --- ``Stigmatisation of people with schizophrenia in Japan''). This itself is a useful target: the \emph{language} of psychiatry slides between describing dopaminergic pathology and accusing the patient of a character defect.
\item
  \textbf{Honest caveat:} I could not locate the exact phrase ``remarkable disturbance of spontaneity'' in the time available. Suggested verification: Kasahara Y. (1984), \emph{Apathy Syndrome} (Iwanami); Shimazono Y. (1974); or possibly an English-language commentary on either. Recommend the author confirm the precise citation before the chapter goes to press; if the phrase is from Bleuler in translation, that is also possible.
\end{itemize}

\textbf{Stimulant withdrawal narratives}

\begin{itemize}
\tightlist
\item
  The clinical picture of post-stimulant anhedonia is substantial. Coffey, Dansky, Carrigan \& Brady cited in PMC3307593 --- anhedonia among cocaine-dependent subjects after 5 days abstinence elevated vs.~controls. \href{https://pmc.ncbi.nlm.nih.gov/articles/PMC2650808/}{PubMed Central} Gawin \& Kleber (1986 \emph{Arch Gen Psychiatry} 43:107--113) --- anhedonia persisted up to 10 weeks post-cessation. \href{https://pmc.ncbi.nlm.nih.gov/articles/PMC2650808/}{PubMed Central} PMC2650808 (Leventhal et al.): in psychiatric outpatients with stimulant use disorder in remission ≥2 months, prevalence of anhedonia and amotivation significantly elevated vs.~controls. \href{https://pmc.ncbi.nlm.nih.gov/articles/PMC2650808/}{PubMed Central} PMC10069411: ``Anhedonia and dysphoria can last for months in people who use MA (D. Hunt et al., 2006; Rawson, 2013).'' \href{https://www.ncbi.nlm.nih.gov/books/NBK576550/}{NCBI}
\item
  For first-person narratives: \emph{David Sheff}, \emph{Beautiful Boy} (2008); \emph{Nic Sheff}, \emph{Tweak} (2008); William Cope Moyers, \emph{Broken} (2006); \emph{Bill Clegg, Portrait of an Addict as a Young Man} (2010). For methamphetamine specifically, \emph{Methland} (Reding 2009) and academic ethnography (Pine, Goffman). For the philosophical/literary register most useful to the manifesto, see Maia Szalavitz, \emph{Unbroken Brain} (2016) --- which weaves first-person ADHD/addiction lived experience with the brain-disease/learning-model debate.
\end{itemize}

\textbf{ADHD lived experience}

\begin{itemize}
\tightlist
\item
  The Aeon essay ``How the hypercuriosity of ADHD may have helped humans thrive'' provides serviceable framing: ``people with ADHD have lower baseline dopamine levels as well as reduced dopamine receptor activity, meaning our brains literally struggle to produce the neurotransmitter necessary to initiate and complete unrewarding tasks.'' \href{https://www.psychologytoday.com/us/blog/mind-the-generation-gap/202601/adhd-and-the-motivation-that-never-comes}{Psychology Today}
\item
  ADHD memoirs to engage: Edward Hallowell \& John Ratey, \emph{Driven to Distraction} (1994) and \emph{Delivered from Distraction} (2005); Sari Solden, \emph{Women with Attention Deficit Disorder} (1995); Lisa Ling, etc. Rebecca Schiller, \emph{A Thousand Ways to Pay Attention} (2021) \href{https://www.goodreads.com/book/show/58328363-a-thousand-ways-to-pay-attention}{Goodreads} is recent and lyrical. Rachel Kramer Bussel's essay ``How an ADHD Memoir Helped Me Stop Feeling Guilty\ldots{}'' is a serviceable first-person passage.
\item
  The phenomenological core that recurs across ADHD memoirs: tasks that are ``important'' but not ``interesting'' produce a substrate-level inability to initiate that is experienced as moral failure.
\end{itemize}

\textbf{Parkinson's first-person}

\begin{itemize}
\tightlist
\item
  Oliver Sacks, \emph{Awakenings} (1973) --- the master text. The post-encephalitic patients (Leonard L., Hester Y., Rose R.) \href{https://grokipedia.com/page/Awakenings_(book)}{Grokipedia} provide the most vivid clinical narratives anywhere of what dopamine collapse and restoration \emph{feel} like. From the book, on Parkinsonian akinesia: ``Parkinsonism is gravity, L-DOPA is levity, and it's difficult to find any mean in between'' \href{http://www.acamedia.info/literature/alexanders_feast/awakenings/excerpts.htm}{Acamedia} (patient quote, p.~10 of 1990 ed.). For the ``gravitational'' sense of effort and struggle the chapter can quote Sacks's reading of Charcot: ``the tremor, rigidity, and akinesia'' \href{http://www.acamedia.info/literature/alexanders_feast/awakenings/excerpts.htm}{Acamedia} as ``the very embodiment'' of resistance to action.
\item
  Michael J. Fox, \emph{Lucky Man} (2002), \emph{Always Looking Up} (2009).
\item
  Jon Palfreman, \emph{Brain Storms} (2015) on PD non-motor symptoms.
\item
  Peter Dunlap-Shohl, \emph{My Degeneration} (graphic memoir, 2015).
\end{itemize}

\textbf{The ``annihilation of the biological precondition for value itself'' --- philosophical anchor}

\begin{itemize}
\tightlist
\item
  The phrase the author wants to develop maps cleanly onto:

  \begin{itemize}
  \tightlist
  \item
    \emph{Bernard Williams} on internal reasons and motivation (1981, ``Internal and External Reasons'' in \emph{Moral Luck}) --- the idea that without the right motivational substrate, reasons fail to bind.
  \item
    \emph{Harry Frankfurt} on the ``wanton'' and on caring (1971, ``Freedom of the Will and the Concept of a Person''; \emph{The Importance of What We Care About}, 1988) --- agency requires second-order desires that \emph{latch}. Dopamine deficit is in this sense a structural production of wantonness.
  \item
    \emph{Susan Wolf}, \emph{Meaning in Life and Why It Matters} (2010) --- meaning requires ``active engagement in projects of worth''; without active engagement (which dopamine subserves) the engagement-side of meaning fails.
  \item
    \emph{Thomas Nagel}, ``The Absurd'' (1971) --- useful counter-text; the chapter can argue \emph{against} Nagel that absurdity isn't just a perspectival fact but can be a substrate fact.
  \item
    \emph{Nietzsche} on ``affective nihilism'' (in \emph{The Will to Power} and \emph{Genealogy of Morals}) --- already a recognized parallel between depressive anhedonia and philosophical nihilism (PsychologyToday 2026 piece, ``Are You a Nihilist or Anhedonic?''). The chapter can run the move productively: what philosophers call nihilism the neurochemist would call striatal hypofunction.
  \end{itemize}
\item
  Pereboom \& Caruso (2018) ``Hard-Incompatibilist Existentialism'' \href{https://cornell.academia.edu/DerkPereboom}{Academia} is the right philosophical anchor for the move from dopamine-determinism to meaning-without-blame.
\end{itemize}

\begin{center}\rule{0.5\linewidth}{0.5pt}\end{center}

\subsection{Suggestions for Vividness and Force}\label{07_pleasure_and_self-suggestions-for-vividness-and-force}

\begin{enumerate}
\def\labelenumi{\arabic{enumi}.}
\tightlist
\item
  \textbf{Open the chapter with an explicit philological-historical move.} Walk the reader from Augustine through Calvin through Olds-Milner through Berridge --- show the cultural arc by which ``pleasure'' was first sin, then suspect science, then linguistically banished from the neuroscience canon. Vrecko's ``civilizing technologies'' thesis is your argumentative pivot.
\item
  \textbf{Use abulia as an unexpected wedge.} Most readers won't know that neurology has, since 1838, had a clinical name --- \emph{non-will} --- for what the chapter is describing. Lead with abulia. The fact that a 19th-century neurology textbook and a 21st-century philosopher of mind are saying the same thing is rhetorically powerful.
\item
  \textbf{Set the wanting/liking critique up by quoting Berridge with respect, then pivoting to humans.} The Weber et al.~2023 amisulpride RCT is your clean experimental hammer. Pair it with the Costello et al.~2025 Parkinson's depression study. Then quote Berridge/Kringelbach 2017 conceding that dopamine is not pleasure ``per se'' --- and ask: per se? On rodent orofacial reactivity? In a human nervous system, ``per se'' is a metaphysical claim, not an empirical one.
\item
  \textbf{Contrast the three cases concretely.} A Parkinson's patient who cannot start his car gets a casserole from the neighbors. An ADHD adult who cannot start her thesis gets called lazy. A meth user who cannot stop using gets arrested. Same striatal circuit. Make this paragraph short and brutal.
\item
  \textbf{Earn the eliminativism.} Don't claim more than necessary. The chapter does not need to argue that \emph{all} of folk psychology is wrong. It needs to argue that the specific term ``willpower'' picks out no biological kind and that, where it is invoked in clinical and moral contexts, it is doing political work, not descriptive work.
\item
  \textbf{Close on Sacks.} \emph{Awakenings} is the book that, more than any other, demonstrates what dopamine \emph{is} by showing what its restoration does and what its loss is. Quote Leonard L.'s diary: ``born again.'' Quote the patient: ``Parkinsonism is gravity, L-DOPA is levity.'' \href{http://www.acamedia.info/literature/alexanders_feast/awakenings/excerpts.htm}{Acamedia} Then make the philosophical claim: this is what we mean by the biological precondition for value.
\end{enumerate}

\begin{center}\rule{0.5\linewidth}{0.5pt}\end{center}

\subsection{Cumulative Arc the Chapter Should Travel}\label{07_pleasure_and_self-cumulative-arc-the-chapter-should-travel}

Section 1 establishes the \emph{cultural} indictment: pleasure was demonized, and that demonization is still alive in neuroscientific terminology. Section 2 establishes the \emph{empirical} indictment: the wanting/liking distinction has been over-extended from rodents to humans, and the human evidence largely reconverges on the older view that dopamine matters for hedonic experience. Section 3 establishes the \emph{philosophical} indictment: ``willpower'' as a folk faculty is at best a poor description of biological computation, and at worst a moralized re-description of dopamine-modulated cortico-striato-thalamo-cortical loops. Section 4 establishes the \emph{political} indictment: same substrate, different moral verdicts --- Parkinson's, ADHD, addiction. Section 5 establishes the \emph{existential} indictment: dopamine system damage is the annihilation of the precondition for value, and the lived experience of those who suffer it is exactly what Bleuler, Kasahara, Sacks, and the abulia tradition have been telling us for 200 years.

\begin{center}\rule{0.5\linewidth}{0.5pt}\end{center}

\subsection{Caveats --- Where Evidence is Mixed or Uncertain}\label{07_pleasure_and_self-caveats-where-evidence-is-mixed-or-uncertain}

\begin{itemize}
\tightlist
\item
  \textbf{Wanting/liking dissociation in humans.} It is \emph{not} nothing. The Brauer \& De Wit 1997 and Leyton 2005/2007 findings are real, \href{https://robinsonlab.research.wesleyan.edu/files/2014/01/Robinson-2015-Curr-Top-Behav-Neurosci.pdf}{Wesleyan} even if their effect sizes are modest. Honest acknowledgment: \emph{some} dissociation can be experimentally produced; the chapter's claim should be that this dissociation does not license the phenomenological claim that dopamine-deficient humans aren't ``really'' suffering hedonically.
\item
  \textbf{The brain-disease model of addiction is contested.} Heather, Lewis, Heilig debate. The chapter can stay agnostic on the specific ``disease'' label while affirming the substrate claim.
\item
  \textbf{Eliminative materialism is a minority position.} Cite Churchland for the family of arguments rather than committing the chapter to full elimination.
\item
  \textbf{The exact Japanese citation for ``remarkable disturbance of spontaneity''} should be verified before publication. Strongest candidates: Kasahara Y. (1984) \emph{Apathy Syndrome} (Iwanami, in Japanese, with English summaries); Shimazono Y. (1974) chapter in Mitsuda \& Fukuda eds.; or possibly a translation of Bleuler. PANSS-G13 (``disturbance of volition'') and N6 (``lack of spontaneity'') are the standardized echoes of the same observation.
\item
  \textbf{Multiple realizability cuts both ways.} Putnam can be cited for the framework but should not be made into an ally of strong reductionism. The careful move is: mental kinds are functional/computational; in \emph{humans} those functions are dopamine-implemented; damage = loss.
\item
  \textbf{Volkow's ``disease of free will'' framing has been cited approvingly here for rhetorical purchase, but it is itself contested.} Some philosophers (Pickard, Heyman) argue it both over-claims compulsion and under-respects patient agency. The chapter's strongest move is to use Volkow's framing as a \emph{concession from inside the establishment}, while not adopting all of its corollaries.
\end{itemize}

\begin{center}\rule{0.5\linewidth}{0.5pt}\end{center}

\subsection{Core References (selected, all real and verifiable)}\label{07_pleasure_and_self-core-references-selected-all-real-and-verifiable}

\textbf{Wanting/liking framework:} - Berridge KC, Robinson TE. Brain Res Rev 1998;28:309--369. - Berridge KC, Kringelbach ML. Neuron 2015;86:646--664. - Berridge KC, Robinson TE. Am Psychol 2016;71:670--679. - Robinson TE, Berridge KC. Annu Rev Psychol 2025;76:29--58.

\textbf{Critiques and effort-based dopamine:} - Salamone JD, Correa M. Neuron 2012;76:470--485. - Salamone JD, Correa M. Annu Rev Psychol 2024;75:1--32. \href{https://www.annualreviews.org/content/journals/10.1146/annurev-psych-020223-012208}{Annual Reviews}

\textbf{Human dopamine-and-pleasure:} - Weber S et al.~J Psychiatr Res 2023 (amisulpride RCT). - Costello H et al.~Brain 2025 (PD depression). - Volkow ND et al.~Mol Psychiatry 2011 (ADHD reward pathway, PMC3010326). - Volkow ND, Wise RA. Nat Neurosci 2005 (obesity/addiction).

\textbf{Olds-Milner and ``pleasure center'':} - Olds J, Milner P. J Comp Physiol Psychol 1954;47:419--427. - Heath RG. J Nerv Ment Dis 1972 (B-19 case).

\textbf{Cultural-history of reward terminology:} - Vrecko S. History of the Human Sciences 2010;23:52--67. - Vrecko S. BioSocieties 2010;5:36--51. - Hammer R et al.~AJOB Neurosci 2013;4:27--32 (PMC3969751).

\textbf{Eliminativism, philosophy of mind, free will:} - Churchland PM. J Philos 1981;78:67--90. - Churchland PS. Neurophilosophy. MIT 1986. - Churchland PS. Touching a Nerve. Norton 2013. - Putnam H. ``Psychological Predicates'' / ``The Nature of Mental States'' 1967. - Pereboom D. Free Will, Agency, and Meaning in Life. OUP 2014. \href{https://philpapers.org/rec/CARPOD-2}{PhilPapers} - Caruso GD. Rejecting Retributivism. CUP 2021. \href{https://www.cambridge.org/core/books/moral-responsibility-reconsidered/0DD200C8E84E6458EA56A098BC7E206F}{Cambridge Core} - Caruso GD, Flanagan O (eds). Neuroexistentialism. OUP 2018.

\textbf{Addiction ethics:} - Heather N. Neuroethics 2017 (PMC5486515). - Pickard H. Neuroethics 2017;10:169--180. - Lewis M. The Biology of Desire. PublicAffairs 2015. - Volkow ND. ``Addiction: A Disease of Free Will.'' NIDA 2015.

\textbf{Anhedonia, abulia, apathy, schizophrenia negative symptoms:} - StatPearls ``Abulia'' (NBK537093). - Foussias G, Remington G. Schizophr Bull 2010;36:359--369 (PMC2833114). - Kasahara Y. Apathy Syndrome. Iwanami 1984 (Japanese). - Furuhashi T. Psychiatry Clin Neurosci Reports 2023 (PMC11114399, on Kasahara).

\textbf{First-person and clinical narrative:} - Sacks O. Awakenings. Duckworth 1973 (rev. 1990). \href{https://grokipedia.com/page/Awakenings_(book)}{Grokipedia} - Fox MJ. Lucky Man. Hyperion 2002. - Szalavitz M. Unbroken Brain. St.~Martin's 2016.
